% spec, 2026-10-24. Section "Exact model and counting". Matches shared/MODEL.md v1.6 §4. Uses \pending from theory's main.tex.
\section{The exact model and the counting}\label{sec:model}

\subsection{What AlphaDev counts}
\emph{Length} in~\cite{alphadev} is the number of x86 instructions of the kernel. Every instruction counts~1: loads
(\texttt{mov} memory$\to$register), stores (\texttt{mov} register$\to$memory), register moves, \texttt{cmp} and
\texttt{cmov}$X$. We checked this on every published kernel: the instruction counts of the listings are exactly the reported
lengths (17, 28, 42, and 43 for the libc++ sort5). The kernels of Table~\ref{tab:lengths} have exactly $n$ loads and $n$
stores (Table~\ref{tab:mix}); some SI variants load a cell twice (e.g.\ \texttt{Sort3\_17}: 4 loads). The data are 32-bit \texttt{int}s in a memory buffer, sorted in place; the conditions
\texttt{l}, \texttt{le}, \texttt{g}, \texttt{ge} are the \emph{signed} ones. Ties are judged by value (the test inputs
contain equal elements).

\paragraph{Instruction set.} $I_0 = \{\texttt{mov}, \texttt{cmp}, \texttt{cmovl}, \texttt{cmovg}\}$ and
$I_1 = I_0 \cup \{\texttt{cmovle}, \texttt{cmovge}\}$. AlphaDev's kernels need $I_1$: the repository sort3 and both sort5
kernels (42, 43) fail if \texttt{le}/\texttt{ge} are replaced by \texttt{l}/\texttt{g}; SI \texttt{Sort3\_17} and both
sort4 kernels still sort after this replacement. A lower bound for AlphaDev must therefore hold in $I_1$. (The shortcut
``on distinct inputs \texttt{cmovle} $=$ \texttt{cmovl}'' is false: a \texttt{cmp} of two registers that hold copies of one
input gives \textsf{eq} also on distinct inputs.) Lengths 17/28/42 are reachable in $I_0$ as well: $I_0$ kernels of these
lengths exist for $n = 3, 4$ (above) and for $n = 5$ (found by our search, \texttt{programs/found/upper\_sort5\_I0\_L42\_mem.s}).

\paragraph{Flags.} \texttt{cmovl} reads $SF \ne OF$, \texttt{cmovle} reads $ZF \lor SF \ne OF$, \texttt{cmovg} reads
$\lnot ZF \land SF = OF$, \texttt{cmovge} reads $SF = OF$. For these four conditions every one of the 16 states of
$(CF, ZF, SF, OF)$ acts as one of three abstract states \textsf{lt}, \textsf{eq}, \textsf{gt}; the state $ZF = 1, SF \ne OF$
acts as \textsf{lt} (Lean theorem \texttt{xeval\_eq}; also checked on hardware).

\subsection{The memory model MM}
\begin{itemize}
\item Memory cells $m_0, \dots, m_{n-1}$ hold the inputs $x_0, \dots, x_{n-1}$ at the start. At the end $m_j$ must hold the
$j$-th smallest input.
\item Registers $r_0, r_1, \dots$: unbounded in number, arbitrary (garbage) values at the start; the flags are arbitrary at the
start.
\item Instructions: \texttt{mov r,[m]} (load), \texttt{mov [m],r} (store), \texttt{mov r,r'}, \texttt{cmp r,r'},
\texttt{cmov}$X$\texttt{ r,r'}, $X$ from the instruction set.
\item \textbf{MM-1}: memory is accessed only by loads and stores. This is AlphaDev's rule ``We cannot compare or conditionally
move to a memory location (illegal)'' (Methods). The x86 ISA also has memory \emph{source} operands (\texttt{cmp r,[m]},
\texttt{cmp [m],r}, \texttt{cmov}$X$\texttt{ r,[m]}); no published AlphaDev kernel uses them. We treat MM-1 as part of the
model and study the variant without it separately (MM+, below).
\item Correct = sorts every $x \in \mathbb{Z}_{32}^n$ for every initial register and flag state. Length = number of
instructions.
\end{itemize}
We do not assume AlphaDev's search heuristics (cells read in order and only once, registers allocated in order, no reads of
uninitialised registers, no \texttt{cmov} before the first \texttt{cmp}, no consecutive \texttt{cmp}s). A lower bound in MM
therefore covers every program AlphaDev's search could produce, and more.

\paragraph{Main claim.} In MM with $I_1$ and unbounded registers, no kernel of length $16$, $27$, $41$ sorts $3$, $4$, $5$
inputs. Equivalently, AlphaDev's lengths $17$, $28$, $42$ are optimal. \pending

\subsection{From memory to registers: the snapshot model}
Search tools work on registers only. The natural register model puts $x_i$ in register $P_i$ and asks for the sorted values
in $P_0, \dots, P_{n-1}$ at the end (model RF, ``fixed outputs''). RF is \emph{too strict}: a store can copy any register to
memory at any time, so AlphaDev's model lets each output come from any register at any step. A lower bound in RF therefore
does not bound MM. The exact register model is the \emph{snapshot model} RS: as RF, plus a free pseudo-instruction
\texttt{out j, r} that records the current value of $r$ as output $j$ (the last \texttt{out j} counts). Its cost $k$ counts
\texttt{mov}, \texttt{cmp} and \texttt{cmov} only.

\begin{lemma}[Snapshot lemma]\label{lem:snapshot}
Under MM-1, with unbounded registers, in the strict and in the relaxed semantics: $L^{\min}_{\mathrm{MM}}(n) = 2n +
k^{\min}_{\mathrm{RS}}(n)$.
\end{lemma}
\noindent\emph{Proof sketch} (full proof: Lemma~L5 in the appendix). ($\le$) Load every $x_i$ into a register first; replace
each \texttt{out j, r} by a store at the same position. ($\ge$) Every cell is loaded before its first store (else $x_i$ is
lost) and stored at least once (else output $j$ is $x_j$, wrong for $n \ge 2$). For each cell, the last load before the first
store is deleted and its target register renamed to $P_i$ for its live range; every store becomes an \texttt{out}, plus a copy
into a fresh register per segment if the cell is read again, and in each such segment the last load is deleted and renamed to
that copy. Per cell this saves one load (before the first store) and one store (the final store, or the last reload after it).
Cost $\le L - 2n$. $\square$

\noindent The construction is implemented (\texttt{agents/spec/tools/ad2out.py}); it maps all 15 published kernels with
$n \le 5$ to RS programs of cost $\le L - 2n$ (sort3: 11, sort4: 20, sort5: 32), and each image passes the strict judge.
Hence the targets in RS are: no program of cost $\le 10$, $\le 19$, $\le 31$ for $n = 3, 4, 5$, in $I_1$.

\paragraph{Fixed outputs at the end.} Let RE be RS with all outputs taken at the end (any registers). Always
$k_{\mathrm{RS}} \le k_{\mathrm{RE}} \le k_{\mathrm{RF}}$. All published kernels still sort when their stores are moved to
the end, so $k_{\mathrm{RE}} \le L_{\mathrm{AD}} - 2n$, and a lower bound $k_{\mathrm{RS}} \ge L_{\mathrm{AD}} - 2n$ gives
$k_{\mathrm{RS}} = k_{\mathrm{RE}}$. For $n = 2$: $k_{\mathrm{RS}} = k_{\mathrm{RE}} = 4$ (exhaustive search without
pruning), so $L^{\min}_{\mathrm{MM}}(2) = 8$; for $n = 3$: $11$. We do not need, and do not claim, $k_{\mathrm{RS}} =
k_{\mathrm{RE}}$ for all $n$.

\subsection{The variant MM+: memory source operands}
MM+ is MM without MM-1: \texttt{cmp r,[m]}, \texttt{cmp [m],r} and \texttt{cmov}$X$\texttt{ r,[m]} are allowed (all are
x86-legal; there is no form with two memory operands and no \texttt{cmov} with a memory destination). Cells are still written
only by stores; every instruction counts~1. The snapshot lemma fails in MM+: a \texttt{cmp} or \texttt{cmov} may read a cell
after its last load (so the load cannot be renamed away), a final store cannot be removed if the cell is read later, and $x_i$
can reach a register by a \texttt{cmov} from memory alone. In general only $L^{\min}_{\mathrm{MM+}} \ge k^{\min}_{\mathrm{RS}}$
holds. MM+ is strictly stronger than MM: for $n = 2$ a 7-instruction kernel exists
(\texttt{mov r0,[m0]; cmp r0,[m1]; mov r1,[m1]; cmovg r1,[m0]; cmovg r0,[m1]; mov [m0],r0; mov [m1],r1}), against 8 in MM;
and our search found MM+ kernels of length 15, 26, 40 for $n = 3, 4, 5$, all checked in the strict judge (sort4 and sort5
also on hardware \pending{} for sort2/sort3).
So AlphaDev's lengths are optimal only \emph{because of} its rule MM-1; under AlphaDev's counting, the full x86 ISA allows
shorter kernels.

\begin{table}[h]\centering\small
\begin{tabular}{lp{7.2cm}p{4.6cm}}\toprule
model & definition & use \\\midrule
MM & memory via loads/stores only (MM-1), all instructions count & main claim (AlphaDev's setting) \\
RS & registers, free snapshot outputs; $L_{\mathrm{MM}} = 2n + k_{\mathrm{RS}}$ & search for the main claim \\
RE, RF & outputs at the end (any / fixed registers) & not a bound for MM \\
MM+ & MM plus memory source operands in \texttt{cmp}, \texttt{cmov} & side result: 15 / 26 / 40 \\\bottomrule
\end{tabular}
\caption{Models. Lower bounds transfer only from RS to MM (Lemma~\ref{lem:snapshot}).}
\end{table}
